\section{Task Construction}
\label{app:tasks}

\subsection{Layered graph generator}

For a prescribed layer-width sequence
$(1,w_1,\ldots,w_L,1)$, vertices receive stable integer identifiers in layer
order. A deterministic coverage construction first adds adjacent-layer edges
so every vertex has route-relevant incoming and outgoing connectivity.
Remaining forward edges, including optional two-layer skips, are sampled
without replacement until the target $m$ is reached. Edges are sorted by
endpoint before their stable indices, costs, and resources are assigned. A
SHA-256-derived seed namespaces the master seed by size stratum and base index,
avoiding dependence on language-level hash randomization.

The v2 strata use layer widths $(1,2,2,1)$,
$(1,2,2,2,1)$, $(1,3,3,1)$, $(1,2,3,3,1)$, and
$(1,3,3,3,1)$ at target edge counts 7, 10, 13, 16, and 19. Five base graphs
are generated per stratum. Positive costs are sampled from $[1,9]$ and
resources from $[1,1000]$ under the frozen master seed 20260827.

\subsection{Route and state enumeration}

All simple source--target routes in a generated DAG are enumerated by a
deterministic traversal. Each route record contains ordered edge identifiers,
objective cost, and resource total. For budget $B$, the exact solver retains
routes with resource total at most $B$ and records all feasible routes tied at
the minimum objective.

Separately, all $2^m$ binary edge selections are enumerated for the benchmark
sizes. For each state, matrix multiplication computes objective cost,
incidence residual, and resource total. A feasible state must have zero flow
residual and meet the budget. This state-domain validator and the path-domain
enumerator agree on the valid route encodings. Their denominators remain
different: $2^m$ for $\phistate$ and the number of candidate routes for
$\phipath$.

\subsection{Duplicate audit and corrected stress schedule}

The historical v1 construction requested seven ordinary quantiles at every
base graph. Of 175 task rows, 53 repeated a feasible set. Ten base graphs had
too few structural path-cardinality levels to supply seven distinct values;
integer resource ties further reduced resolution. The audit retained these
rows as historical evidence rather than modifying them.

For v2, let $R_1<\cdots<R_J$ be the distinct route-resource thresholds and
let $n_j$ be the number of routes with resource exactly $R_j$. The achievable
feasible counts are cumulative sums $N_j=\sum_{k\leq j}n_k$. Preferred counts
are reserved when exactly achievable. An unavailable preferred count $u$ is
mapped to the unused $N_j$ minimizing $|\log(N_j/u)|$, with smaller cardinality
as deterministic tie-breaker. Remaining capacity is filled by the unused
count maximizing its minimum log-distance from selected counts. The sorted
selected thresholds become D1, D2, and so forth; the construction emits at
most seven levels and never repeats a feasible set.

The wider resource support yields 169 distinct route-resource totals across
169 route entries and zero resource-tie graphs. Ten graphs still have fewer
than seven achievable sets and are labeled as insufficient for seven levels;
their available distinct levels are retained without duplication. The final
task counts and fractions are shown in \cref{tab:task-strata}.


\begin{table}[ht]
\centering
\caption{Distinct-cardinality task strata used to define the controlled universe.}
\label{tab:task-strata}
\small
\begin{tabular}{@{}lrrrrr@{}}
\toprule
Stratum & Edges & Graphs & Tasks & $\phi_{\min}$ & $\phi_{\max}$ \\
\midrule
S1 & 7 & 5 & 15 & 0.00781 & 0.0234 \\
S2 & 10 & 5 & 20 & 0.000977 & 0.00391 \\
S3 & 13 & 5 & 35 & 0.000122 & 0.000854 \\
S4 & 16 & 5 & 35 & 1.53e-05 & 0.000153 \\
S5 & 19 & 5 & 35 & 1.91e-06 & 2.29e-05 \\
\bottomrule
\end{tabular}
\end{table}


\subsection{Construction scope}

The controlled schedule deliberately increases within-graph resolution. It
does not assert that the preferred cardinalities occur naturally, that all
graphs have comparable combinatorial structure, or that $D$ alone captures
their difficulty. Task generation and selection never use a QAOA result.
